% Preámbulo común del carrusel y de la historia (estilo de los apuntes)
\usepackage{fontspec}
\usepackage[spanish,es-noshorthands]{babel}
\usepackage{microtype}
\usepackage{amsmath, amssymb, mathtools, bm}
\usepackage{xcolor}
\definecolor{fondo}{HTML}{000000}
\definecolor{texto}{HTML}{ECECEC}
\definecolor{blanco}{HTML}{FFFFFF}
\definecolor{negro}{HTML}{000000}
\definecolor{gris}{HTML}{9CA0A8}
\definecolor{grismedio}{HTML}{5B5F67}
\definecolor{grisoscuro}{HTML}{2E3036}
\definecolor{cajafondo}{HTML}{0E1014}
\definecolor{cian}{HTML}{20B4CE}
\definecolor{cianclaro}{HTML}{7FDBEA}
\definecolor{violeta}{HTML}{A48BFA}
\definecolor{ambar}{HTML}{F4B740}
\definecolor{magenta}{HTML}{E5579D}
\definecolor{verde}{HTML}{4FD68A}
\definecolor{naranja}{HTML}{FA8F45}
\definecolor{rojo}{HTML}{F25F5C}
\pagecolor{fondo}\color{texto}
\usepackage{graphicx}
\usepackage{tikz}
\usetikzlibrary{arrows.meta, calc, positioning}
\usepackage{pgfplots}
\pgfplotsset{compat=1.18}
\usepgfplotslibrary{fillbetween}
% =========================================================================
%  Estilos gráficos comunes (TikZ / pgfplots) para fondo negro
% =========================================================================

% --- Mapa de colores «arcoíris» para coloreado por fase (dominio complejo) ---
% Recorre el círculo cromático completo: el color codifica arg f(z) en (-pi,pi].
\pgfplotsset{colormap={fase}{
  rgb255=(255,40,60) rgb255=(255,150,20) rgb255=(250,230,40)
  rgb255=(60,220,90) rgb255=(30,200,230) rgb255=(60,90,255)
  rgb255=(190,60,240) rgb255=(255,40,60)}}

% --- Estilos TikZ generales ------------------------------------------------
\tikzset{
  >=Stealth,
  every picture/.append style={line cap=round, line join=round},
  eje/.style={-{Stealth[length=5pt,width=4pt]}, line width=0.6pt, gris},
  ejeclaro/.style={-{Stealth[length=5pt,width=4pt]}, line width=0.7pt, texto},
  fino/.style={line width=0.4pt, grismedio},
  guia/.style={densely dashed, line width=0.45pt, grismedio},
  punteado/.style={densely dotted, line width=0.6pt, gris},
  etiq/.style={font=\small, text=texto},
  etiqgris/.style={font=\footnotesize, text=gris},
  polo/.style={cross out, draw=rojo, line width=1.1pt, minimum size=6pt, inner sep=0pt},
  punto/.style={circle, fill=blanco, inner sep=0pt, minimum size=4pt},
  puntocian/.style={circle, fill=cian, inner sep=0pt, minimum size=4.5pt},
  orienta/.style={decoration={markings, mark=at position #1 with
     {\arrow{Stealth[length=6pt,width=5pt]}}}, postaction={decorate}},
  orienta/.default=0.55,
  nodo/.style={draw=grismedio, fill=cajafondo, rounded corners=2pt,
     line width=0.5pt, inner xsep=6pt, inner ysep=4pt, font=\small, text=texto},
  nodocian/.style={nodo, draw=cian},
  nodoambar/.style={nodo, draw=ambar},
  nodovioleta/.style={nodo, draw=violeta},
  flecha/.style={-{Stealth[length=5pt,width=4.5pt]}, line width=0.6pt, gris},
}

% --- Estilo base de ejes pgfplots en oscuro ----------------------------------
\pgfplotsset{
  oscuro/.style={
    axis background/.style={fill=fondo},
    tick label style={font=\scriptsize, text=gris},
    label style={font=\small, text=texto},
    title style={font=\small, text=texto},
    axis line style={grismedio, line width=0.5pt},
    every tick/.style={grismedio},
    major tick length=3pt,
    grid=major, grid style={grisoscuro, line width=0.35pt},
    legend style={font=\scriptsize, draw=grisoscuro, fill=cajafondo,
                  text=texto, rounded corners=1.5pt, inner sep=3pt,
                  row sep=1pt},
    legend cell align=left,
    /pgf/number format/use comma,
    /pgf/number format/1000 sep={},
    unbounded coords=jump,
  },
  % Ejes cruzados por el origen (estilo «pizarra»)
  cruz/.style={oscuro, axis lines=middle, grid=none,
    axis line style={gris, line width=0.6pt,
      -{Stealth[length=5pt,width=4pt]}},
    xlabel style={anchor=west, font=\small},
    ylabel style={anchor=south, font=\small},
  },
  superficie/.style={
    tick label style={font=\scriptsize, text=gris},
    label style={font=\small, text=texto},
    axis line style={grismedio}, every tick/.style={grismedio},
    grid=none, colormap name=fase, shader=interp,
    /pgf/number format/use comma,
  },
  % Curvas
  c1/.style={line width=1.15pt, cian},
  c2/.style={line width=1.15pt, magenta},
  c3/.style={line width=1.15pt, ambar},
  c4/.style={line width=1.15pt, verde},
  c5/.style={line width=1.15pt, violeta},
  c6/.style={line width=1.15pt, naranja},
}


% Escala progresiva para familias de curvas (azul -> violeta)
\definecolor{cm1}{RGB}{60,90,255}
\definecolor{cm2}{RGB}{45,145,240}
\definecolor{cm3}{RGB}{30,200,230}
\definecolor{cm4}{RGB}{60,220,90}
\definecolor{cm5}{RGB}{250,230,40}
\definecolor{cm6}{RGB}{255,150,20}
\definecolor{cm7}{RGB}{255,40,60}
\definecolor{cm8}{RGB}{190,60,240}

\pagestyle{empty}
\setlength{\parindent}{0pt}

\DeclareMathOperator{\sen}{sen}
\renewcommand{\sin}{\sen}
\newcommand{\Lop}{\mathcal{L}}
\newcommand{\dif}{\mathrm{d}}
\newcommand{\eul}{\mathrm{e}}
\newcommand{\T}{\mathsf{T}}
\newcommand{\yv}{\hat{\bm y}}
\newcommand{\ubr}[2]{\underbrace{#1}_{\textcolor{gris}{#2}}}
\newcommand{\R}{\mathbb{R}}

% Caja con pestaña, en el estilo de los apuntes
\tikzset{caja/.style={draw=#1, fill=cajafondo, rounded corners=2.5pt,
    line width=0.7pt, inner xsep=8pt, inner ysep=7pt, text width=8.6cm,
    align=flush center, text=texto},
  pestana/.style={fill=#1, rounded corners=2pt, inner xsep=5pt, inner ysep=1.6pt,
    font=\sffamily\bfseries\scriptsize, text=negro},
  chip/.style={draw=#1, fill=cajafondo, rounded corners=2.5pt, line width=0.7pt,
    inner xsep=6pt, inner ysep=6pt, minimum width=4.3cm, minimum height=1.05cm,
    align=flush center, text=texto}}

